\centerline{\bf \Huge Вспоминаем комбу II}
\centerline{\bf \Huge Базированный граф}

\begin{flushright}
        \scriptsize{Мама, вызывай такси!}
\end{flushright}

\problem{1}
$\bigl[\textbf{С какого-то там матбоя}\bigr]$
В особой больнице есть главврач и много пациентов. В течение недели каждый пациент один раз в день кусал кого-нибудь (возможно и себя). В конце недели оказалось, что у всех пациентов по 2 укуса, а у главврача - 100 укусов. Сколько в больнице пациентов?

\problem{2}
Каждый мальчик знаком с 7 мальчиками и 10 девочками, а каждая девочка - с 8 девочками и 9 мальчиками. Кого больше - мальчиков или девочек?

\problem{3}
В углах доски $3 \times 3$ стоят кони: по 2 коня черного и белого цветов, причем кони одного цвета стоят в противоположных углах. За один ход разрешается выбрать любого коня и сделать им ход в свободную клетку. Можно ли переставить коней так, чтобы по прежнему все кони стояли в углах, но кони одного цвета стояли в углах при одной стороне?


\problem{4}
На доске $9 \times 9$ угловые клетки черные. Может ли хромой чернопольный слон (он может ходить только на соседнюю клетку по диагонали) посетить все чёрные клетки, если нельзя проходить через одну клетку дважды?

\problem{5}
После нескольких игровых дней однокругового футбольного чемпионата выяснилось, что любые пять команд можно так расположить по кругу, чтобы каждая команда сыграла со стоящими справа и слева. Докажите, что чемпионат можно завершить в три дня (в один день команда может сыграть не более одной игры).

\problem{6}
Степени всех вершин графа не меньше 3. Докажите, что в нем существует хотя бы один чётный цикл.

\problem{7} Нарисован выпуклый многоугольник, разбитый непересекающимися диагоналями на треугольники. Можно ли стороны и диагонали раскрасить в жёлтый и красный цвета так, чтобы жук мог проползти из каждой вершины в любую другую по жёлтым отрезкам, а клоп - по красным?

\problem{8} Двадцать городов соединены 172 авиалиниями. Докажите, что, используя эти авиалинии, можно из любого города перелететь в любой другой (быть может, делая пересадки).

\newpage

\hspace{3mm}


\problem{9} В стране 15 городов, некоторые из них соединены авиалиниями, принадлежащими трём авиакомпаниям. Известно, что даже если любая из авиакомпаний прекратит полеты, можно будет добраться из каждого города в любой другой (возможно, с пересадками), пользуясь рейсами оставшихся двух компаний. Какое наименьшее количество авиалиний может быть в стране?